% GENERATED from content/data/enrichment.json by mxbuild.enrichment; do not edit.
\mxunit{enrichment-bayes-exercises}{Inverser une probabilité conditionnelle}
\mxsubtitle{Première--Terminale · probabilités \textperiodcentered{} Exercices progressifs \textperiodcentered{} 07}

\begin{note}
\textbf{Données.} Pour les exercices 1 et 2 : on choisit l’urne A avec probabilité $\frac{1}{3}$, l’urne B avec probabilité $\frac{2}{3}$. A contient 3 boules rouges et 1 bleue ; B contient 1 rouge et 3 bleues. On tire ensuite une boule au hasard dans l’urne choisie.
\end{note}

\begin{exercise}
Dans l’exemple, calculer $P(A \cap  \text{bleue})$.
\end{exercise}

\begin{exercise}
Calculer $P(A|\text{bleue})$.
\end{exercise}

\begin{exercise}
Dans une classe, 12 élèves sur 30 jouent aux échecs ; parmi eux 9 aiment les énigmes. Calculer $P(\text{énigmes}|\text{échecs})$. Peut-on en déduire $P(\text{échecs}|\text{énigmes})$ ?
\end{exercise}

\begin{exercise}
Deux événements A et R vérifient $P(A)=\num{0,4}$, $P(R|A)=\num{0,5}$ et $P(R|\text{non }A)=\num{0,5}$. Calculer $P(A|R)$.
\end{exercise}

